% GENERATED FILE. Edit canonical lessons, then run npm run generate:books.
% canonical-source-hash: fnv1a64:041fd3b519f8b865
% canonical-lessons: TA-C102-fatherselderbrother, TA-C102-periyamma, TA-C102-motherinlaw, TA-C102-daughterinlaw, TA-C102-soninlaw

\chapter{Elders and In-Laws}
\label{ch:ta-elders-and-in-laws}
\hlchaptermodality{102}

\begin{chapteropening}
\textbf{By the end of this chapter:} \emph{I can name my father's elder brother and his wife, a mother-in-law, and a daughter- or son-in-law.}

Complete the last lesson of chapter 102: i can name my father's elder brother and his wife, a mother-in-law, and a daughter- or son-in-law.
\end{chapteropening}

\section[periyappā]{\ta{பெரியப்பா} (periyappā) — the father's elder brother}

\label{lesson:TA-C102-fatherselderbrother}



\begin{quote}
Before the new one: say the Tamil for an aunt, then the Tamil for the father's younger brother.
\end{quote}

\subsection*{You'll want to know: \ta{பெரியப்பா}}
\textbf{\ta{பெரியப்பா}} (\emph{periyappā}) — your father's elder brother, the \emph{big father}.

\textbf{\ta{பெரிய}} is \textquotedblleft{}big\textquotedblright{}. Beside \textbf{\ta{சித்தப்பா}}, the little father.

\begin{grammarlens}[title={What you've built}]
Big father and little father.
\end{grammarlens}

\subsection*{Your turn}
\begin{itemize}
\raggedright
  \item \emph{Say it:} \emph{periyappā}
  \item \emph{Say it:} \emph{periyappā}, once more
  \item \emph{Say it:} say \emph{eṉ periyappā}
  \item \emph{Recall:} say the Tamil for an aunt, then the Tamil for the father's younger brother, then say \emph{periyappā} again
\end{itemize}

\begin{tcolorbox}[breakable,colback=teal!4,colframe=teal!35!black,title={Before you move on}]
What does \ta{பெரியப்பா} mean? (\textquotedblleft{}The father's elder brother\textquotedblright{}.) Say it once more.
\end{tcolorbox}
\section[periyammā]{\ta{பெரியம்மா} (periyammā) — the big mother (an elder aunt)}

\label{lesson:TA-C102-periyamma}



\begin{quote}
Before the new one: say the Tamil for the father's younger brother, then the Tamil for the father's elder brother.
\end{quote}

\subsection*{You'll want to know: \ta{பெரியம்மா}}
\textbf{\ta{பெரியம்மா}} (\emph{periyammā}) — the \emph{big mother}: your mother's elder sister, or the wife of your \textbf{\ta{பெரியப்பா}}.

\textbf{\ta{அம்மா}} is the mother you know.

\begin{grammarlens}[title={What you've built}]
The big mother.
\end{grammarlens}

\subsection*{Your turn}
\begin{itemize}
\raggedright
  \item \emph{Say it:} \emph{periyammā}
  \item \emph{Say it:} \emph{periyammā}, once more
  \item \emph{Say it:} say \emph{eṉ periyammā}
  \item \emph{Recall:} say the Tamil for the father's younger brother, then the Tamil for the father's elder brother, then say \emph{periyammā} again
\end{itemize}

\begin{tcolorbox}[breakable,colback=teal!4,colframe=teal!35!black,title={Before you move on}]
What does \ta{பெரியம்மா} mean? (\textquotedblleft{}The big mother (an elder aunt)\textquotedblright{}.) Say it once more.
\end{tcolorbox}
\section[māmiyār]{\ta{மாமியார்} (māmiyār) — a mother-in-law}

\label{lesson:TA-C102-motherinlaw}



\begin{quote}
Before the new one: say the Tamil for the father's elder brother, then the Tamil for the big mother.
\end{quote}

\subsection*{You'll want to know: \ta{மாமியார்}}
\textbf{\ta{மாமியார்}} (\emph{māmiyār}) — \textquotedblleft{}a mother-in-law\textquotedblright{}, the respectful word.

It is built on the same \textbf{\ta{மாமா}} root as the uncle.

\begin{grammarlens}[title={What you've built}]
An in-law.
\end{grammarlens}

\subsection*{Your turn}
\begin{itemize}
\raggedright
  \item \emph{Say it:} \emph{māmiyār}
  \item \emph{Say it:} \emph{māmiyār}, once more
  \item \emph{Say it:} say \emph{eṉ māmiyār}
  \item \emph{Recall:} say the Tamil for the father's elder brother, then the Tamil for the big mother, then say \emph{māmiyār} again
\end{itemize}

\begin{tcolorbox}[breakable,colback=teal!4,colframe=teal!35!black,title={Before you move on}]
What does \ta{மாமியார்} mean? (\textquotedblleft{}A mother-in-law\textquotedblright{}.) Say it once more.
\end{tcolorbox}
\section[marumakaḷ]{\ta{மருமகள்} (marumakaḷ) — a daughter-in-law}

\label{lesson:TA-C102-daughterinlaw}



\begin{quote}
Before the new one: say the Tamil for the big mother, then the Tamil for a mother-in-law.
\end{quote}

\subsection*{You'll want to know: \ta{மருமகள்}}
\textbf{\ta{மருமகள்}} (\emph{marumakaḷ}) — \textquotedblleft{}a daughter-in-law\textquotedblright{}.

\textbf{\ta{மகள்}} is \textquotedblleft{}daughter\textquotedblright{}; \textbf{\ta{மரு}-} marks the relation by marriage.

\begin{grammarlens}[title={What you've built}]
A daughter by marriage.
\end{grammarlens}

\subsection*{Your turn}
\begin{itemize}
\raggedright
  \item \emph{Say it:} \emph{marumakaḷ}
  \item \emph{Say it:} \emph{marumakaḷ}, once more
  \item \emph{Say it:} say \emph{eṉ marumakaḷ}
  \item \emph{Recall:} say the Tamil for the big mother, then the Tamil for a mother-in-law, then say \emph{marumakaḷ} again
\end{itemize}

\begin{tcolorbox}[breakable,colback=teal!4,colframe=teal!35!black,title={Before you move on}]
What does \ta{மருமகள்} mean? (\textquotedblleft{}A daughter-in-law\textquotedblright{}.) Say it once more.
\end{tcolorbox}
\section[marumakaṉ]{\ta{மருமகன்} (marumakaṉ) — a son-in-law}

\label{lesson:TA-C102-soninlaw}



\begin{quote}
Before the new one: say the Tamil for a mother-in-law, then the Tamil for a daughter-in-law.
\end{quote}

\subsection*{You'll want to know: \ta{மருமகன்}}
\textbf{\ta{மருமகன்}} (\emph{marumakaṉ}) — \textquotedblleft{}a son-in-law\textquotedblright{}.

\textbf{\ta{மகன்}} is \textquotedblleft{}son\textquotedblright{}; the same \textbf{\ta{மரு}-} makes him family by marriage.

\begin{grammarlens}[title={What you've built}]
Elders and in-laws.
\end{grammarlens}

\subsection*{Your turn}
\begin{itemize}
\raggedright
  \item \emph{Say it:} \emph{marumakaṉ}
  \item \emph{Say it:} \emph{marumakaṉ}, once more
  \item \emph{Say it:} say \emph{eṉ marumakaṉ}
  \item \emph{Recall:} say the Tamil for a mother-in-law, then the Tamil for a daughter-in-law, then say \emph{marumakaṉ} again
\end{itemize}

\begin{tcolorbox}[breakable,colback=teal!4,colframe=teal!35!black,title={Before you move on}]
What does \ta{மருமகன்} mean? (\textquotedblleft{}A son-in-law\textquotedblright{}.) Say it once more.
\end{tcolorbox}
